\documentclass{article}

% Preprint draft setup
\usepackage[utf8]{inputenc}
\usepackage[english]{babel}
\usepackage{amsmath,amssymb,amsfonts}
\usepackage{graphicx}
\usepackage{hyperref}
\usepackage{booktabs}
\usepackage{algorithm}
\usepackage{algorithmic}
\usepackage[margin=1in]{geometry}

% Title and authors
\title{HumanBrain: GPU-Accelerated Whole-Brain Simulator with Anatomical Connectivity and Adaptive Feedback Control}

\author{
    Francisco Molina Burgos \\
    \texttt{pako.molina@gmail.com} \\
    ORCID: 0009-0008-6093-8267
}

\date{November 2025}

\begin{document}

\maketitle

\begin{abstract}
We present HumanBrain, a GPU-accelerated computational framework for whole-brain simulation exploring multi-compartmental cable equation physics, literature-referenced anatomical connectivity, and adaptive feedback control. The framework implements: (1) wgpu compute shaders for 152-compartment cable equation dynamics on GPU (1.52M compartments total for 10K neurons), (2) eight literature-derived inter-regional pathways including thalamocortical, corticothalamic, corticostriatal, and hippocampal-cortical connections, and (3) a hybrid CPU-GPU feedback loop designed for attractor analysis and homeostatic modulation. Exploratory single-run profiles observe frame rates of ~50-80 FPS on consumer hardware (NVIDIA RTX 3050; formal reproducible benchmark suite pending). Implemented submodules provide structural prototypes for hippocampal place cells, basal ganglia action selection, and thalamic burst/tonic firing modes (unverified against experimental recordings). HumanBrain provides an open framework for computational neuroscience research.
\end{abstract}

\section{Introduction}

Whole-brain simulation remains a grand challenge in computational neuroscience.

\section*{Code Availability}

Source code: \url{https://github.com/Yatrogenesis/HumanBrain}

\end{document}
